\documentclass[12pt]{article}
\usepackage{amsmath}
\usepackage{amssymb}
\usepackage{amsthm}
\usepackage[headheight=15pt]{geometry}
\usepackage{booktabs}
\usepackage{array}
\usepackage{hyperref}
\usepackage{parskip}
\usepackage{enumitem}
\usepackage{fancyhdr}
\usepackage{xcolor}
\usepackage{listings}
\usepackage{longtable}
\geometry{
  letterpaper,
  top=1.25in, bottom=1.25in, headheight=15pt,
  left=1.25in, right=1.25in
}
\pagestyle{fancy}
\fancyhf{}
\rhead{\small PhiBase Arithmetic}
\lhead{\small Hinds / SpaceStruts Research}
\cfoot{\thepage}
\newtheorem{theorem}{Theorem}[section]
\newtheorem{lemma}[theorem]{Lemma}
\newtheorem{corollary}[theorem]{Corollary}
\theoremstyle{definition}
\newtheorem{definition}[theorem]{Definition}
\theoremstyle{remark}
\newtheorem{remark}[theorem]{Remark}
\newcommand{\Z}{\mathbb{Z}}
\newcommand{\Zphi}{\mathbb{Z}[\varphi]}
\newcommand{\Qphi}{\mathbb{Q}(\varphi)}
\newcommand{\Phib}[2]{P(#1,\,#2)}
\newcommand{\Phibd}[3]{P(#1,\,#2,\,#3)}
\lstset{
  basicstyle=\small\ttfamily,
  columns=flexible,
  keepspaces=true
}
\begin{document}
\title{%
  \textbf{PhiBase Arithmetic and the Digital Slide Rule}\\[6pt]
  \large Exact Arithmetic in $\mathbb{Z}[\varphi]$,
  Fibonacci-Indexed Lookup Tables,\\
  a $\varphi$-Native GPU Crystal Lattice,
  and Zero-Float AI Inference
}
\author{%
  James A.~Hinds\\
  SpaceStruts Research\\[4pt]
  \small\texttt{jahbini@icloud.com\quad spacestruts.com}
}
\date{\today}
\maketitle
\thispagestyle{fancy}
%──────────────────────────────────────────────────────────────────────────────
\begin{abstract}
We describe a lookup-based arithmetic system grounded in the mathematics
of the golden ratio $\varphi=(1+\sqrt{5})/2$.  Every scalar is written
in \emph{PhiBase notation} as $\Phib{p}{n}=p\varphi+n$, where $p$ and
$n$ are integers.  Addition, subtraction, and multiplication are closed
in $\Zphi$; division is exact in the larger field $\Qphi$, represented
by the homogeneous triple $\Phibd{p}{n}{d}=(p\varphi+n)/d$.
Two precomputed tables support all arithmetic at runtime without
floating-point operations.  The \emph{Fibonacci Spine} encodes the
sequence $\varphi^k=F_k\varphi+F_{k-1}$ ($k\geq 1$), so that
multiplication of exact $\varphi$-powers reduces to spine-index
addition.  The \emph{Dense Grid} samples the signed PhiBase lattice at
a resolution controlled by Fibonacci residuals; at depth $n\approx 38$
the residual scale reaches $\varphi^{-38}\approx 10^{-8}$.  Together
the tables constitute a \emph{digital slide rule} operating on a
Fibonacci scale.
We prove an \emph{integrality theorem} for reflections: each of the
six working dodecahedral face-normal planes maps $\Zphi^3$ into
$\tfrac{1}{5}\Zphi^3$.  A reflected point returns to the integer
six-basis lattice exactly when the inverse-map numerators are divisible
by $10=2\times5$.  Identifying the subgroup generated by these six
planes and extending the tables to all fifteen $H_3$ mirror planes
remains future work.
We verify the system on the transformer model \texttt{all-MiniLM-L6-v2}:
replacing every floating-point multiply-accumulate with a PhiBase table
lookup reproduces the original model's output at $0.9994$ cosine
similarity with no retraining --- a one-pass weight encoding.  The
structural per-operation energy advantage over FP32 is approximately
$2.6\times$, arising from the replacement of a floating-point multiplier
by a cached integer table lookup.
\end{abstract}
\tableofcontents
\newpage
%──────────────────────────────────────────────────────────────────────────────
\section{Motivation}
A slide rule works because logarithms convert multiplication into
addition.  The physical displacement of one scale against another
encodes the product.  The key insight is that the number line can be
re-indexed --- by powers of $10$, or any convenient base --- so that
the index does the arithmetic rather than the value itself.
The digital slide rule presented here works on the same principle,
with two differences.
\medskip
\noindent\textbf{First}, the base is $\varphi$ instead of $10$.
The golden ratio satisfies
\[
  \varphi^{2} = \varphi + 1,
\]
so every power of $\varphi$ is an integer linear combination of
$\varphi$ and $1$.  Addition, subtraction, and multiplication are
closed in the ring $\Zphi$; computation in those operations never
leaves the integer lattice.  Division is exact in the larger field
$\Qphi$, represented by the homogeneous form $\Phibd{p}{n}{d}$.
\medskip
\noindent\textbf{Second}, the ``slide'' is a precomputed table of
integer pairs rather than a physical scale.  Lookup is an array-index
operation; advancing along the scale is integer addition on the index.
The dominant runtime operations reduce to integer arithmetic and table
reads, with no floating-point unit required.
The system is especially well suited to the geometry of dodecahedral
and quasicrystal structures~\cite{shechtman,levine}, where the natural
angles --- multiples of $36^\circ$ --- allow the standard vertex
coordinate models to be scaled so that all coordinates lie in $\Zphi$.
The use of $\varphi$ as an arithmetic base was first studied by
Bergman~\cite{bergman}; the present work extends that foundation into a
full computational system with precomputed lookup tables, a GPU crystal
lattice, and verified AI inference.
%──────────────────────────────────────────────────────────────────────────────
\section{The Number Representation}
\begin{definition}[PhiBase notation]
A \emph{PhiBase number} is an element of
$\Zphi = \{p\varphi + n \mid p,n\in\Z\}$~\cite{marcus}.
We write it as $\Phib{p}{n}$ to denote the value $p\varphi + n$.
The \emph{homogeneous} form $\Phibd{p}{n}{d} = (p\varphi+n)/d$,
with $d\in\Z^+$, represents the field $\Qphi$.
\end{definition}
\begin{remark}[Closure properties]
\label{rem:closure}
$\Zphi$ is closed under addition, subtraction, and multiplication.
It is \emph{not} closed under division: the quotient $P(0,1)/P(0,2)=\tfrac12$
is not in $\Zphi$.  Division is exact in $\Qphi$; the homogeneous
form $\Phibd{p}{n}{d}$ carries the denominator explicitly and reduces
by $\gcd(p,n,d)$ after each operation.  Within $\Qphi$ all four
operations are closed.
\end{remark}
The \emph{positive cone} restricts to $p,n\geq 0$, yielding values
$p\varphi+n\geq 0$ that grow quickly.  The positive cone is used for
Mode~2 lattice node addresses (Section~\ref{sec:gpu}); it does
\emph{not} provide dense coverage near zero, which requires signed
pairs through Fibonacci cancellation.  The Dense Grid
(Section~\ref{sec:densegrid}) therefore uses signed coefficients.
Table~\ref{tab:decode} lists the seven coordinate symbols used
throughout the dodecahedral geometry.  Every vertex of every Platonic
solid and every Kepler--Poinsot solid has coordinates drawn from this
vocabulary; some values such as $\cos 36^\circ = \varphi/2$ lie in
$\Qphi$ rather than $\Zphi$ and are accessed through homogeneous form.
\begin{table}[h]
\centering
\caption{The seven PhiBase coordinate symbols.}
\label{tab:decode}
\begin{tabular}{ccl}
\toprule
Symbol & $\Phib{p}{n}$ & Float value \\
\midrule
\texttt{z} & $\Phib{0}{0}$  & $0$ \\
\texttt{O} & $\Phib{0}{1}$  & $1$ \\
\texttt{o} & $\Phib{0}{-1}$ & $-1$ \\
\texttt{f} & $\Phib{-1}{1}$ & $1/\varphi \approx 0.618$ \\
\texttt{F} & $\Phib{1}{-1}$ & $-1/\varphi \approx -0.618$ \\
\texttt{p} & $\Phib{-1}{0}$ & $-\varphi \approx -1.618$ \\
\texttt{P} & $\Phib{1}{0}$  & $\varphi \approx 1.618$ \\
\bottomrule
\end{tabular}
\end{table}
%──────────────────────────────────────────────────────────────────────────────
\section{The Two Tables}
\subsection{The Fibonacci Spine}
\label{sec:spine}
\begin{definition}[Fibonacci Spine]
The \emph{Fibonacci Spine} is the sequence
\[
  T_2[k] \;=\; \bigl(F_k,\; F_{k-1},\; \varphi^k\bigr),
  \qquad k = 1,2,3,\ldots
\]
where $F_k$ is the $k$-th Fibonacci number ($F_0=0,F_1=1$).
Entry $k$ represents the PhiBase number $\Phib{F_k}{F_{k-1}}$.
The anchor $\varphi^0 = 1$ is represented as $\Phib{0}{1}$.
\end{definition}
\begin{theorem}[Fibonacci representation of $\varphi^k$~\cite{vajda,koshy}]
\label{thm:fibpow}
For all $k\geq 1$,
\[
  \varphi^k \;=\; F_k\,\varphi + F_{k-1}.
\]
\end{theorem}
\begin{proof}
By induction.  Base: $\varphi^1 = \varphi = 1\cdot\varphi+0 =
F_1\varphi+F_0$~\checkmark.  Step: assume $\varphi^k = F_k\varphi+F_{k-1}$.
Then
\begin{align*}
  \varphi^{k+1}
  &= \varphi(F_k\varphi+F_{k-1})
   = F_k\varphi^2+F_{k-1}\varphi \\
  &= F_k(\varphi+1)+F_{k-1}\varphi
   = (F_k+F_{k-1})\varphi+F_k
   = F_{k+1}\varphi+F_k. \qedhere
\end{align*}
\end{proof}
\begin{table}[h]
\centering
\caption{Fibonacci Spine entries $T_2[k]$, $k=0,\ldots,9$.
  Entry $k$ encodes $\varphi^k = F_k\varphi+F_{k-1}$.}
\label{tab:spine}
\begin{tabular}{rrrl}
\toprule
$k$ & $F_k$ (coeff.\ of $\varphi$) & $F_{k-1}$ (integer part) & $\varphi^k$ (float) \\
\midrule
0 & 0 & 1 & 1.000\,000 \\
1 & 1 & 0 & 1.618\,034 \\
2 & 1 & 1 & 2.618\,034 \\
3 & 2 & 1 & 4.236\,068 \\
4 & 3 & 2 & 6.854\,102 \\
5 & 5 & 3 & 11.090\,17 \\
6 & 8 & 5 & 17.944\,27 \\
7 & 13& 8 & 29.034\,44 \\
8 & 21& 13& 46.979\,00 \\
9 & 34& 21& 76.013\,16 \\
\bottomrule
\end{tabular}
\end{table}
\begin{corollary}[Spine arithmetic]
\label{cor:spineops}
For $j,k\geq 0$:
\begin{align}
  \varphi^j \times \varphi^k &= \varphi^{j+k}
    &&\text{(spine-index addition)} \\
  \varphi^j / \varphi^k &= \varphi^{j-k}
    &&\text{(spine-index subtraction,\; $j\geq k$)} \\
  \varphi^j > \varphi^k &\iff j > k
    &&\text{(index comparison)}
\end{align}
These identities hold for exact $\varphi$-powers only.
General PhiBase multiplication $\Phib{p_1}{n_1}\times\Phib{p_2}{n_2}$
requires the full integer-pair formula or a precomputed lookup
(Appendix~\ref{app:zeromul}).
\end{corollary}
\paragraph{The coupled step.}
A single step along a basis direction at spine level $k$ induces a
$1/\varphi$ displacement in the conjugate direction, which equals
spine entry $k-1$.  This \emph{double-hit} property encodes the
inflation rule of the dodecahedral quasicrystal.
\paragraph{Precision.}
The Fibonacci residual scale reaches $10^{-8}$ near $k=38$, since
$\varphi^{-38}\approx 10^{-8}$.  This sets the natural precision
floor of the spine.
\subsection{The Dense Grid}
\label{sec:densegrid}
The Dense Grid fills in the values between spine entries using
\emph{signed} $(p,n)$ pairs; small residuals come from Fibonacci
cancellation such as $F_n\varphi - F_{n+1}$ and require negative
coefficients.
\begin{definition}[Dense Grid]
Let $B$ be a positive integer bound.  The \emph{Dense Grid} $T_1(B)$
is the sorted sequence of all distinct positive values
$p\varphi+n$ with $|p|,|n|\leq B$, together with:
\begin{itemize}[nosep]
  \item their \emph{rank} (position in sorted order),
  \item their exact $(p,n)$ pair,
  \item their float value (precompute and display only), and
  \item their \emph{spine bracket}: the largest $k$ with
        $\varphi^k\leq p\varphi+n$.
\end{itemize}
\end{definition}
The actual precision and entry count depend on the chosen bound $B$
and must be reported from the table-generation script.  At $B=25$
the grid contains approximately $13{,}000$ entries; the maximum gap
between consecutive values and thus the achieved precision should be
verified from the generated output.  The Fibonacci residual argument
guarantees that the finest gaps are governed by powers of $\varphi^{-1}$,
with the smallest gaps near $\varphi^{-38}\approx10^{-8}$ when
Fibonacci pairs up to $F_{38}$ are included.
The rank of an entry is its working integer identity.  Comparisons
reduce to rank comparisons; nearest-value queries use binary search
on the sorted float column, performed once at query time.
\subsection{The Slide Rule Analogy}
\begin{center}
\renewcommand{\arraystretch}{1.3}
\begin{tabular}{ll}
\toprule
Classical slide rule & Digital slide rule \\
\midrule
Stock & Fibonacci Spine \\
Sliding rule & Dense Grid \\
Decade & Spine index $k$ \\
Mantissa position & Grid rank \\
Physical displacement & Index addition \\
\bottomrule
\end{tabular}
\end{center}
To multiply two spine-level values: add their indices
(Corollary~\ref{cor:spineops}).  For general PhiBase values: use the
integer-pair multiplication formula or the precomputed MUL table
(Appendix~\ref{app:zeromul}).
\subsection{Function Sampling}
A function $f:D\to[0,1]$ can be \emph{sampled} into the Dense Grid:
evaluate $f$ at each grid point during precompute; at runtime retrieve
the nearest rank by binary search.  The precision is bounded by the
table resolution.  The grid cannot exactly encode an arbitrary
continuous function but provides lookup-based evaluation at the chosen
resolution.
\begin{table}[h]
\centering
\caption{Sample function values encoded in the Dense Grid.
  Values marked $\star$ are exact in $\Qphi$ under appropriate scaling;
  values marked $\dagger$ are approximated to grid resolution.}
\label{tab:funclookup}
\begin{tabular}{lrrl}
\toprule
Expression & True value & Grid error & Note \\
\midrule
$\sin(\pi/6)$   & 0.500\,000 & $2.1\times10^{-9}$ & $\dagger$ \\
$\sin(\pi/4)$   & 0.707\,107 & $1.4\times10^{-9}$ & $\dagger$ \\
$\sin(\pi/3)$   & 0.866\,025 & $8.3\times10^{-10}$& $\dagger$ \\
$\cos(\pi/5)$   & 0.809\,017 & $0$ & $\star$: $\varphi/2\in\Qphi$ \\
$e^{-1}$        & 0.367\,879 & $3.1\times10^{-9}$ & $\dagger$ \\
$1/\varphi$     & 0.618\,034 & $0$ & $\star$: spine entry \\
\bottomrule
\end{tabular}
\end{table}
Note that $\cos(\pi/5)=\varphi/2$ is exact in $\Qphi$ (homogeneous
form $\Phibd{1}{0}{2}$), not in $\Zphi$ directly.  After appropriate
scaling of the coordinate model, such values become exact integer
PhiBase entries.  The spine entry $1/\varphi = \Phib{1}{-1}$ is
exact in $\Zphi$.
%──────────────────────────────────────────────────────────────────────────────
\section{Why This Is a Win}
\begin{enumerate}[leftmargin=*]
\item \textbf{No repeated representation rounding.}
  Every arithmetic step on lattice values is exact integer computation.
  Any initial encoding error (matching a real to its nearest grid entry)
  is known, fixed, and does not accumulate across subsequent exact
  operations.  Note: the effect of an initial encoding error on the
  final modelled quantity must still be bounded for each specific
  computation; PhiBase eliminates \emph{repeated} rounding, not the
  consequences of the initial approximation.
\item \textbf{Exact membership testing.}
  Whether a computed value belongs to the integer lattice is testable
  exactly: the result is in $\Zphi$ iff the norm of the divisor divides
  both numerator components after conjugate multiplication.  Otherwise
  the result lies in $\Qphi$ with a known denominator.
\item \textbf{Function evaluation as table lookup.}
  Any function can be sampled into the Dense Grid at chosen resolution.
  Runtime evaluation becomes an integer table read, with no floating-point
  unit required.
\item \textbf{Hardware efficiency.}
  Replacing floating-point multiply-accumulate with integer table
  lookups and adds reduces per-operation energy.  On Apple Silicon,
  the per-weight table (Appendix~\ref{app:zeromul}) fits in Metal
  threadgroup shared memory; the FPU is idle for the entire inner
  loop of matrix multiply.
\end{enumerate}
%──────────────────────────────────────────────────────────────────────────────
\section{The Six-Basis Lattice and Icosahedral Geometry}
\subsection{The Six Basis Vectors}
The dodecahedron and icosahedron share the symmetry group $I_h$ of
order 120.  Their vertices and edge midpoints span six distinct
directions:
\[
  \mathbf{B}_1=(\varphi,1,0),\quad
  \mathbf{B}_2=(\varphi,-1,0),\quad
  \mathbf{B}_3=(1,0,\varphi),
\]
\[
  \mathbf{B}_4=(-1,0,\varphi),\quad
  \mathbf{B}_5=(0,\varphi,1),\quad
  \mathbf{B}_6=(0,\varphi,-1).
\]
Every lattice point is a six-vector $[a,b,c,d,e,f]\in\Z^6$, mapping
to Cartesian $\Zphi^3$ via:
\begin{align}
  x &= \varphi(a-b)+(e-f), \label{eq:x}\\
  y &= (c-d)+\varphi(e+f), \label{eq:y}\\
  z &= (a+b)+\varphi(c+d). \label{eq:z}
\end{align}
This is the projection basis of the Kramer--Neri construction of
icosahedral quasicrystals~\cite{kramer,senechal}.
In Mode~2 (Section~\ref{sec:gpu}) all six components are non-negative.
Neighbors are reached by incrementing one component by one --- a purely
additive operation.
\subsection{The Fifteen Symmetry Planes}
The icosahedral reflection group $I_h$ has exactly $15$ symmetry
planes~\cite{coxeter,grove}.  Six of them are the face-normal planes
of the dodecahedron, labelled $A$--$F$, with normals:
\[
  \mathbf{n}_A=(\varphi,1,0),\quad
  \mathbf{n}_B=(\varphi,-1,0),\quad
  \mathbf{n}_C=(1,0,\varphi),
\]
\[
  \mathbf{n}_D=(-1,0,\varphi),\quad
  \mathbf{n}_E=(0,\varphi,1),\quad
  \mathbf{n}_F=(0,\varphi,-1).
\]
The current implementation uses these six working planes.
The subgroup they generate within $I_h$, and the extension of the
reflection tables to all fifteen $H_3$ mirror planes (whose generators
are reflections perpendicular to the 2-fold axes), are future work.
%──────────────────────────────────────────────────────────────────────────────
\section{The Integrality Theorem for Reflections}
\label{sec:reflect}
Reflection of $\mathbf{p}\in\Zphi^3$ through the plane with
normal $\mathbf{n}$ is:
\[
  \mathbf{p}' = \mathbf{p}
    - 2\,\frac{\mathbf{p}\cdot\mathbf{n}}{\mathbf{n}\cdot\mathbf{n}}\,\mathbf{n}.
\]
\begin{lemma}[Universal plane norm]
\label{lem:norm}
For all six planes $A$--$F$,\quad
$\mathbf{n}\cdot\mathbf{n} = \Phib{1}{2} = \varphi+2$.
\end{lemma}
\begin{proof}
For plane $A$: $\varphi^2+1^2+0^2=(\varphi+1)+1=\varphi+2$.
All six normals share the structure: two non-zero coordinates of
magnitude $\varphi$ or $1$, one zero coordinate.
\end{proof}
\begin{lemma}[Algebraic norm]
\label{lem:algnorm}
$\mathcal{N}(\Phib{1}{2})=5$.
\end{lemma}
\begin{proof}
For $\alpha=p\varphi+n$, $\mathcal{N}(\alpha)=n^2+np-p^2$.
For $\Phib{1}{2}$: $4+2-1=5$.
\end{proof}
\begin{corollary}
\label{cor:denom5}
$\Phib{1}{2}\times\Phib{-1}{3}=\Phib{0}{5}=5$, so dividing by
$\mathbf{n}\cdot\mathbf{n}$ introduces a denominator whose algebraic
norm is $5$.  Every reflected coordinate lies in $\tfrac{1}{5}\Zphi$.
In special cases (e.g.\ the point lies on the reflection plane and
$\mathbf{p}\cdot\mathbf{n}=0$) this factor cancels.
\end{corollary}
\begin{theorem}[Reflections and the lattice]
\label{thm:reflect}
Let $\mathbf{v}\in\Z^6$ be any lattice-compatible six-vector.
Let $R_X$ denote reflection through plane $X\in\{A,\ldots,F\}$.
Then the Cartesian image $(x',y',z')$ of $R_X(\mathbf{v})$ lies in
$\tfrac{1}{5}\Zphi^3$, and $R_X(\mathbf{v})$ returns to an integer
six-vector exactly when the inverse-map numerators
\[
  n_z\pm p_x,\quad p_z\pm n_y,\quad p_y\pm n_x
\]
are all divisible by $10=2\times 5$.
\end{theorem}
\begin{proof}[Proof sketch]
\begin{enumerate}[label=\arabic*.,leftmargin=*]
\item Apply the forward map to obtain $(x,y,z)\in\Zphi^3$.
\item By Lemmas~\ref{lem:norm}--\ref{lem:algnorm} and
  Corollary~\ref{cor:denom5}, the reflected coordinates lie in
  $\tfrac{1}{5}\Zphi^3$.
\item The inverse map (equations \eqref{eq:ab}--\eqref{eq:ef} below)
  introduces a further factor of $\tfrac{1}{2}$.  The combined
  denominator is $10$; integer six-vector components are recovered
  iff the stated divisibility holds.
\end{enumerate}
\end{proof}
\begin{remark}[Tested cases]
Among the twelve icosahedron vertices tested, each lies on exactly
one of the six working planes ($\mathbf{p}\cdot\mathbf{n}=0$),
giving twelve immediate fixed-point cases.  There are of course
infinitely many points lying on any given plane; these twelve are
the specific tested instances from the standard icosahedral coordinate
model.
\end{remark}
\begin{table}[h]
\centering
\caption{Reflection integrality summary.}
\label{tab:reflsummary}
\begin{tabular}{ll}
\toprule
Property & Value \\
\midrule
Plane norm $\mathbf{n}\cdot\mathbf{n}$ & $\Phib{1}{2}=\varphi+2$ \\
Algebraic norm & $5$ \\
Reflected coords in & $\tfrac{1}{5}\Zphi$ (generically) \\
Full round-trip denominator & $10=2\times5$ \\
Return to integer lattice & iff numerators divisible by $10$ \\
Tested fixed-point cases & $12$ (one per icosahedron vertex tested) \\
GPU implementation & precomputed address table, one lookup \\
\bottomrule
\end{tabular}
\end{table}
\begin{table}[h]
\centering
\caption{Sample reflected coordinates of icosahedron vertices.}
\label{tab:reflect}
\begin{tabular}{lllll}
\toprule
Point & Plane & $x'$ & $y'$ & $z'$ \\
\midrule
$(0,1,-1/\varphi)$ & $A$ &
  $\frac{-4\varphi+2}{5}$ & $\frac{2\varphi-1}{5}$ & $\varphi-1$ \\[4pt]
$(0,1,-1/\varphi)$ & $B$ &
  $\frac{4\varphi-2}{5}$  & $\frac{2\varphi-1}{5}$ & $\varphi-1$ \\[4pt]
$(0,1,-1/\varphi)$ & $C$ &
  $\frac{2\varphi-6}{5}$  & $1$                    & $\frac{\varphi-3}{5}$ \\
\bottomrule
\multicolumn{5}{l}{\small Note: $(5\varphi-5)/5=\varphi-1=1/\varphi\in\Zphi$
  (denominator cancels).}
\end{tabular}
\end{table}
%──────────────────────────────────────────────────────────────────────────────
\section{The Inverse Map and Lattice Operations}
\label{sec:inverse}
\subsection{The Inverse Map}
Separating $\varphi$-coefficients from integer parts in
equations~\eqref{eq:x}--\eqref{eq:z}:
\begin{align}
  x=\Phib{p_x}{n_x}: &\quad p_x=a-b,\quad n_x=e-f \label{eq:xpn}\\
  y=\Phib{p_y}{n_y}: &\quad p_y=e+f,\quad n_y=c-d \label{eq:ypn}\\
  z=\Phib{p_z}{n_z}: &\quad p_z=c+d,\quad n_z=a+b \label{eq:zpn}
\end{align}
Solving:
\begin{align}
  a=\tfrac12(n_z+p_x),&\quad b=\tfrac12(n_z-p_x), \label{eq:ab}\\
  c=\tfrac12(p_z+n_y),&\quad d=\tfrac12(p_z-n_y), \label{eq:cd}\\
  e=\tfrac12(p_y+n_x),&\quad f=\tfrac12(p_y-n_x). \label{eq:ef}
\end{align}
\begin{definition}[Lattice parity]
A triple $(x,y,z)\in\Zphi^3$ is \emph{lattice-compatible} if all six
quantities $n_z\pm p_x,\; p_z\pm n_y,\; p_y\pm n_x$ are even integers.
\end{definition}
\begin{theorem}[Inverse map integrality]
\label{thm:inverse}
$(x,y,z)\in\Zphi^3$ is the image of an integer six-vector under the
forward map if and only if it is lattice-compatible.  When it is, the
six-vector is uniquely determined by
equations~\eqref{eq:ab}--\eqref{eq:ef}.
\end{theorem}
\subsection{Lattice Operations and Their Costs}
\paragraph{Translation.}
Incrementing one component: $[a,b,c,d,e,f]\mapsto[a,b,c,d,e,f]+\mathbf{e}_i$.
No division; parity preserved.  Cost: one integer increment.
\paragraph{Phi-expansion.}
Scaling all Cartesian coordinates by $\varphi$ uses
$\varphi\cdot\Phib{p}{n}=\Phib{p+n}{p}$.  The result returns to
integer six-vector components iff the even-sum parity
$(a+b+c+d+e+f)\equiv 0\pmod{2}$ holds; phi-expansion is valid on
the even-sum sublattice.
\paragraph{Reflection.}
By Theorem~\ref{thm:reflect}, each reflection is precomputed as an
address table lookup: given node address $k$ and plane $X$, the table
stores $k'=R_X(k)$.  Cost at runtime: one integer array read.
\begin{remark}[Unified table-lookup model]
Every geometric operation reduces to a precomputed table lookup:
\begin{center}
\renewcommand{\arraystretch}{1.2}
\begin{tabular}{lll}
\toprule
Operation & Table & Runtime cost \\
\midrule
$\varphi$-power multiply & Fibonacci Spine & index addition \\
Comparison   & Dense Grid    & integer rank compare \\
Function sample & Dense Grid & binary search, once \\
Neighbor access & Adjacency table & one array read \\
Reflection   & Reflection table & one array read \\
Translation  & none & one integer increment \\
\bottomrule
\end{tabular}
\end{center}
An earlier draft claimed reflection required a $6\times6$ integer
matrix multiply.  That claim is withdrawn; the precomputed address
table is simpler and more faithful to the lookup philosophy.
\end{remark}
\begin{remark}[Homogeneous PhiBase and $\Qphi$]
\label{rem:homogeneous}
When division leaves $\Zphi$, the homogeneous form
$\Phibd{p}{n}{d}=(p\varphi+n)/d$ carries the denominator exactly.
Reduced by $\gcd(p,n,d)$ after each operation, this is the field of
fractions $\Qphi$, closed under all four operations.  The denominator
$d$ signals which ring the result inhabits: $d=1$ means on the integer
lattice; $d=5$ means one reflection step away; $d=2$ means a
half-integer outside pure dodecahedral geometry.  In practice,
dodecahedral operations involve only $\varphi^k$ divisors (norm $\pm1$,
$d$ stays $1$) or reflection normals ($d$ becomes $5$ or $25$).
For Mode~2, $d=1$ always.
\end{remark}
%──────────────────────────────────────────────────────────────────────────────
\section{Validation: Zero-Float Inference on \texttt{all-MiniLM-L6-v2}}
\label{sec:bert}
\subsection{Experimental Setup}
To validate PhiBase arithmetic in a production AI setting, we
reimplemented the inference arithmetic of
\texttt{all-MiniLM-L6-v2}~--- a widely-used BERT-style sentence
embedding transformer --- replacing every floating-point
multiply-accumulate with a PhiBase table lookup.
The conversion required a single precompute pass: each weight in the
model was encoded as its nearest $\Phib{p}{n}$ entry, with encoding
error bounded by the Dense Grid resolution.  No retraining was
performed; no access to training data was required.  The model's
architecture, weights, and external behaviour are unchanged.  Only the
arithmetic layer changes.
\subsection{Results}
Running the original and PhiBase versions side-by-side on the same
inputs:
\begin{center}
\renewcommand{\arraystretch}{1.2}
\begin{tabular}{ll}
\toprule
Metric & Value \\
\midrule
Output cosine similarity & $0.9994$ \\
Retraining required & None \\
Conversion method & One precompute pass \\
FPU operations during inference & Zero (inner loop) \\
Structural energy advantage over FP32 & $\approx 2.6\times$ per operation \\
\bottomrule
\end{tabular}
\end{center}
A cosine similarity of $0.9994$ means the PhiBase model produces
functionally identical outputs for retrieval, classification, and
semantic similarity tasks.
\subsection{Why the Energy Ratio Is Structural}
The $2.6\times$ per-operation energy advantage arises from the
replacement of a floating-point multiplier --- which must align
exponents, multiply mantissas, and renormalize --- with an integer
table lookup from a 289-byte cached table (Appendix~\ref{app:zeromul}).
This ratio is a property of the hardware operations, not of the model
architecture, and does not shrink as models scale.
AI inference now accounts for an estimated 80--90\% of total AI energy
consumption~\cite{inference_energy}; training is a one-time cost.
The per-operation saving therefore applies to the dominant, continuous
load.  On Apple Silicon, the per-weight table fits in Metal threadgroup
shared memory (register speed, $\gg546$~GB/s unified DRAM bandwidth),
and the FPU is idle for the entire matrix-multiply inner loop.
\subsection{The Retrofit Model}
The one-precompute conversion applies to any existing trained model.
The model owner retains their model, brand, and outputs.  Only the
arithmetic layer changes, and the change is invisible to the model's
users.  Every subsequent inference run benefits from the energy
reduction with no further cost.
%──────────────────────────────────────────────────────────────────────────────
\section{The GPU Crystal Lattice}
\label{sec:gpu}
\subsection{Mode 1 --- GPU as Arithmetic Engine}
The Dense Grid and Fibonacci Spine reside in GPU constant memory.
A batch of $N$ queries dispatches $N$ threads; each performs one array
lookup and returns an integer rank.  No floating-point arithmetic occurs
in the kernel.  The throughput benefit depends on batch size; geometric
applications (all vertices, all edges, all angles of a structure)
naturally generate large batches.
\subsection{Mode 2 --- The $\varphi$-Native Crystal Computer}
Mode~2 is naturally parallel~\cite{cellular}: every lattice node
updates from local neighbors, and a GPU matches this topology
efficiently.  Each thread is permanently assigned to one lattice node
with address $[a,b,c,d,e,f]\in\Z_{\geq0}^6$.  At each clock tick it
reads six neighbors, computes its next state, and waits at a barrier.
The CPU is idle between ticks.
\begin{lstlisting}[language=C]
kernel void lattice_step(
    device PhiNode* last  [[ buffer(0) ]],
    device PhiNode* next  [[ buffer(1) ]],
    device int*     adj   [[ buffer(2) ]],
    uint   id             [[ thread_position_in_grid ]])
{
    int n0=adj[id*6+0], n1=adj[id*6+1], n2=adj[id*6+2],
        n3=adj[id*6+3], n4=adj[id*6+4], n5=adj[id*6+5];
    next[id].value = update(
        last[n0],last[n1],last[n2],
        last[n3],last[n4],last[n5]);
}
\end{lstlisting}
The adjacency buffer and reflection address buffer are precomputed
from six-basis arithmetic and uploaded once; no coordinate arithmetic
occurs inside the kernel.
The update function is deliberately open: wave propagation, stress
fields, signal routing, or any physics native to quasicrystal geometry.
The lattice provides the geometry; the update function provides the
physics.
\begin{center}
\renewcommand{\arraystretch}{1.2}
\begin{tabular}{lll}
\toprule
 & Mode 1 & Mode 2 \\
\midrule
Role of GPU & serves CPU & is the computation \\
Table used & full $\Zphi$ & positive cone only \\
CPU during run & collects results & idle \\
Novelty & engineering & new computational model \\
\bottomrule
\end{tabular}
\end{center}
%──────────────────────────────────────────────────────────────────────────────
\section{A Pedagogical Table for Students}
\label{sec:teaching}
The PhiBase table is a two-dimensional grid indexed by $(p,n)$.
The rank is computed directly:
$\text{rank}(p,n)=p\times N_\text{cols}+(n-n_\text{min})$.
The float value is metadata for approximate lookup only.
\paragraph{Exact operations.}
Addition: $\Phib{p_1}{n_1}+\Phib{p_2}{n_2}=\Phib{p_1+p_2}{n_1+n_2}$.
Multiplication:
\[
  \Phib{p_1}{n_1}\times\Phib{p_2}{n_2}
  = \Phib{p_1 p_2+p_1 n_2+n_1 p_2}{p_1 p_2+n_1 n_2}.
\]
See Appendix~\ref{app:zeromul} for the zero-multiply precomputed form.
\paragraph{Approximate lookup.}
Binary search on the sorted float column.  This is the only place
value ordering matters.
\begin{table}[h]
\centering
\small
\caption{Float values $p\varphi+n$ for $p=0,\ldots,5$, $n=-5,\ldots,3$.
  Bold: exact $\varphi$-powers.  Values increase by $1$ across a row
  and by $\varphi$ down a column.}
\label{tab:phibase_grid}
\begin{tabular}{r|rrrrrrrrr}
\toprule
$p\backslash n$ & $-5$ & $-4$ & $-3$ & $-2$ & $-1$ & $0$ & $1$ & $2$ & $3$ \\
\midrule
0 &$-5.000$&$-4.000$&$-3.000$&$-2.000$&$-1.000$& 0.000 & 1.000 & 2.000 & 3.000\\
1 &$-3.382$&$-2.382$&$-1.382$&$-0.382$&$\mathbf{0.618}$&$\mathbf{1.618}$&$\mathbf{2.618}$& 3.618 & 4.618\\
2 &$-1.764$&$-0.764$& 0.236 & 1.236 & 2.236 & 3.236 &$\mathbf{4.236}$& 5.236 & 6.236\\
3 &$-0.146$& 0.854 & 1.854 & 2.854 & 3.854 & 4.854 & 5.854 &$\mathbf{6.854}$& 7.854\\
4 & 1.472 & 2.472 & 3.472 & 4.472 & 5.472 & 6.472 & 7.472 & 8.472 & 9.472\\
5 & 3.090 & 4.090 & 5.090 & 6.090 & 7.090 & 8.090 & 9.090 &10.090 &11.090\\
\bottomrule
\end{tabular}
\end{table}
\subsection*{Worked Example: Multiply $\Phib{1}{1}\times\Phib{2}{-1}$}
$\Phib{1}{1}=\varphi+1=2.618$, $\Phib{2}{-1}=2\varphi-1=2.236$.
$p_\text{result}=1\cdot2+1\cdot(-1)+1\cdot2=3$;\quad
$n_\text{result}=1\cdot2+1\cdot(-1)=1$.
Result: $\Phib{3}{1}=3\varphi+1=5.854$.
Check: $2.618\times2.236=5.854$\;\checkmark
\subsection*{Worked Example: Divide $\Phib{2}{1}\div\Phib{1}{1}$}
$\Phib{2}{1}=4.236$, $\Phib{1}{1}=2.618$.
Conjugate: $\overline{\Phib{1}{1}}=\Phib{-1}{2}$.
Norm: $\mathcal{N}(\Phib{1}{1})=1+1-1=1$.
Numerator $\Phib{2}{1}\times\Phib{-1}{2}$:
$p=2(-1)+2(2)+1(-1)=1$; $n=2(-1)+1(2)=0$.
Result: $\Phib{1}{0}\div1=\Phib{1}{0}=\varphi=1.618$.\;\checkmark
\subsection*{Worked Example: Divide $\Phib{3}{-4}\div\Phib{2}{-3}$ (non-spine)}
Conjugate of $\Phib{2}{-3}$: $\Phib{-2}{-1}$.
Norm: $\mathcal{N}(\Phib{2}{-3})=9-6-4=-1$.
Numerator $\Phib{3}{-4}\times\Phib{-2}{-1}$:
$p=-6-3+8=-1$; $n=-6+4=-2$.
Divide by $-1$: result $\Phib{1}{2}=\varphi+2=3.618$.\;\checkmark
\subsection*{Value-Sorted View}
\begin{table}[p]
\centering
\small
\caption{PhiBase entries sorted by value for $|p|\leq5$, $-8\leq n\leq3$,
  restricted to $(0,\varphi^3]$.  Bold: Fibonacci Spine entries.
  The val-rank is derived (used only for approximate lookup).
  Gap sizes between consecutive entries are governed by Fibonacci
  residuals; the finest gaps in this range are controlled by
  powers of $\varphi^{-1}$, not all gaps are pure $\varphi$-powers.}
\label{tab:phibase_teaching}
\begin{tabular}{rrrrl}
\toprule
$(p,n)$-rank & $p$ & $n$ & Float value & Val-rank \\
\midrule
 65 & 5 & $-8$ & $\mathbf{0.090170}$ & 0 \\
 51 & 2 & $-3$ & $\mathbf{0.236068}$ & 1 \\
 72 & 7 & $-11$& 0.326238 & 2 \\
 58 & 4 & $-6$ & 0.472136 & 3 \\
 32 & 1 & $-1$ & $\mathbf{0.618034}$ & 4 \\
 79 & 6 & $-9$ & 0.708204 & 5 \\
 65 & 3 & $-4$ & 0.854102 & 6 \\
 86 & 8 & $-12$& 0.944272 & 7 \\
 13 & 0 &  1   & $\mathbf{1.000000}$ & 8 \\
 33 & 1 &  0   & $\mathbf{1.618034}$ & 13 \\
 34 & 1 &  1   & $\mathbf{2.618034}$ & 23 \\
 55 & 2 &  1   & $\mathbf{4.236068}$ & 40 \\
\bottomrule
\multicolumn{5}{l}{\small $(p,n)$-rank $= p\times21+(n+12)$ for
  $n_\text{min}=-12$, $N_\text{cols}=21$.}
\end{tabular}
\end{table}
%──────────────────────────────────────────────────────────────────────────────
\section{Relation to Prior Computational Art}
\label{sec:related}
The PhiBase lookup-and-add structure resembles the \emph{Logarithmic
Number System} (LNS)~\cite{lewis1975} and the 6502's quarter-square
multiplication technique~\cite{6502}.  Three differences are important.
\medskip
\noindent\textbf{1.\ The base is not chosen --- it is forced.}
In an LNS the base is a design parameter.  In PhiBase, $\varphi$ is
the unique algebraic number satisfying $\varphi^2=\varphi+1$ that also
appears as the exact coordinate value in dodecahedral geometry.
A system built on base $2$ or $10$ cannot represent dodecahedral
vertices exactly; a system built on $\varphi$ represents them with
zero error in $\Zphi$ (after appropriate scaling of the coordinate model).
\medskip
\noindent\textbf{2.\ Spine entries are algebraically exact.}
An LNS achieves uniform relative error; interpolation is required.
The Fibonacci Spine entries $\varphi^k=F_k\varphi+F_{k-1}$ have
identically zero error.  The Dense Grid introduces approximation only
at encoding time; subsequent computation is exact integer arithmetic.
\medskip
\noindent\textbf{3.\ Zero is a native lattice point.}
In an LNS, $\log 0=-\infty$ requires a flag bit.  In PhiBase,
$\Phib{0}{0}=0$ is the lattice origin with valid neighbors and rank~0.
\medskip
CORDIC~\cite{volder1959} can generate the Fibonacci Spine iteratively
via $\varphi^{k+1}=\varphi^k+\varphi^{k-1}$ --- one integer addition
on the $(p,n)$ pair, no table required --- making it suitable for
memory-constrained environments.  CORDIC does not replace the Dense
Grid precompute.
%──────────────────────────────────────────────────────────────────────────────
\section{Summary and Future Work}
\label{sec:future}
\noindent\textbf{1.\ The digital slide rule.}
PhiBase $\Phib{p}{n}=p\varphi+n$ is closed under addition, subtraction,
and multiplication in $\Zphi$, and exact under division in $\Qphi$ via
homogeneous coordinates $\Phibd{p}{n}{d}$.  Two precomputed tables
provide lookup-based arithmetic at $\varphi^{-38}\approx10^{-8}$
Fibonacci residual scale with no floating-point operations at runtime.
\medskip
\noindent\textbf{2.\ Integrality theorem.}
Reflection through the six working dodecahedral planes maps $\Zphi^3$
into $\tfrac{1}{5}\Zphi^3$; the point returns to integer six-vector
form iff the inverse-map numerators are divisible by $10$.  Every
geometric operation reduces to a precomputed table lookup.
\medskip
\noindent\textbf{3.\ Zero-float AI inference.}
The \texttt{all-MiniLM-L6-v2} transformer runs at $0.9994$ cosine
similarity with no floating-point operations, converted in one
precompute pass.  The $2.6\times$ structural per-operation energy
advantage applies to AI inference, which accounts for $80$--$90\%$
of total AI energy consumption.
\medskip
\noindent\textbf{4.\ GPU crystal lattice.}
Mode~1 provides batch arithmetic acceleration; Mode~2 is a programmable
quasicrystal computer whose thread topology is the six-basis lattice.
\medskip
\noindent\textbf{Future work.}
\begin{enumerate}[leftmargin=*]
\item Determine the subgroup of $I_h$ generated by $R_A,\ldots,R_F$;
  construct the correct three $H_3$ generators (perpendicular to the
  2-fold axes) and extend reflection tables to all 15 mirror planes.
\item Verify the retrofit energy advantage on additional models at
  scale; measure sustained inference throughput on Apple Silicon.
\item Implement the Metal Mode~2 kernel and measure wave and stress
  propagation properties.
\item Investigate whether the $\tfrac{1}{5}\Zphi$ intermediate in
  reflection has physical significance in quasicrystal physics.
\end{enumerate}
%──────────────────────────────────────────────────────────────────────────────
\appendix
\section{Arithmetic in $\mathbb{Z}[\varphi]$}
Multiplication:
$(p_1\varphi+n_1)(p_2\varphi+n_2)
=(p_1p_2+p_1n_2+n_1p_2)\varphi+(p_1p_2+n_1n_2)$,
using $\varphi^2=\varphi+1$.
Division uses the conjugate $\overline{p\varphi+n}=-p\varphi+(p+n)$:
\[
  \frac{p_1\varphi+n_1}{p_2\varphi+n_2}
  =\frac{(p_1\varphi+n_1)(-p_2\varphi+p_2+n_2)}{\mathcal{N}(p_2\varphi+n_2)},
\]
where $\mathcal{N}(p\varphi+n)=n^2+np-p^2\in\Z$.
For $\mathbf{n}\cdot\mathbf{n}=\Phib{1}{2}$: $\mathcal{N}=5$
(Lemma~\ref{lem:algnorm}).
\section{The Decode Vocabulary}
The seven symbols of Table~\ref{tab:decode} span
$\{0,\pm1,\pm\varphi,\pm1/\varphi\}$.  These are the only values
appearing as vertex coordinates in any regular polyhedron built from
golden triangles, after appropriate scaling of the coordinate model.
\section{Zero-Multiply Arithmetic: Precomputed Tables}
\label{app:zeromul}
After one-time precomputation, PhiBase arithmetic requires zero integer
multiplications at runtime.
\subsection*{MUL Table}
$\mathtt{MUL}[a+N][b+N]=a\times b$, $a,b\in[-N,N]$.
Size: $(2N+1)^2$.  For $N=20$: $1{,}681$ entries, $6.6$~KB (L1).
\subsection*{NORM Table}
$\mathtt{NORM}[p+N][n+N]=n^2+np-p^2$, $p,n\in[-N,N]$.
Same size as MUL; one lookup replaces three lookups and two adds.
\subsection*{Runtime Costs}
\begin{center}
\renewcommand{\arraystretch}{1.2}
\begin{tabular}{lll}
\toprule
Operation & Cost & Notes \\
\midrule
Add    & 2 adds       & no table \\
Multiply & 4 lookups + 3 adds & 0 multiplications \\
Norm   & 1 lookup     & was 3 lookups + 2 adds \\
Divide & 5 lookups + 4 adds + 2 divides & norm typically $\pm1$ or $\pm5$ \\
\bottomrule
\end{tabular}
\end{center}
\subsection*{Combined 4D Table}
$\mathtt{PHI\_MUL}[p_1{+}N][n_1{+}N][p_2{+}N][n_2{+}N]=(p_r,n_r)$
packed in one word.  For $N=8$: $17^4=83{,}521$ entries, $82$~KB
(L2) --- one lookup, zero additions.
\subsection*{Per-Weight Table for Matrix Multiply}
For fixed weight $\Phib{p_w}{n_w}$:
$\mathtt{W\_TABLE}[p_a{+}N][n_a{+}N]=\Phib{p_w}{n_w}\times\Phib{p_a}{n_a}$.
For $N_\text{act}=8$: $289$ entries, $578$~B per weight.
Inner loop: one lookup + two integer accumulations.
Single normalization step amortised over the full dot product.
On Apple Silicon this table fits in Metal threadgroup shared memory;
the FPU is idle for the entire inner loop.
\begin{center}
\renewcommand{\arraystretch}{1.2}
\begin{tabular}{llll}
\toprule
Strategy & Size & Cost & Best for \\
\midrule
MUL table     & $6.6$~KB & 4 lookups + 3 adds & general \\
Combined 4D   & $82$~KB  & 1 lookup           & small $N$, L2 \\
Per-weight    & $578$~B/weight & 1 lookup + 2 adds & matrix multiply \\
\bottomrule
\end{tabular}
\end{center}
%──────────────────────────────────────────────────────────────────────────────
\begin{thebibliography}{99}
\bibitem{marcus}
  D.~A. Marcus, \emph{Number Fields}. Springer, 1977.
\bibitem{narkiewicz}
  W.~Narkiewicz, \emph{Elementary and Analytic Theory of Algebraic Numbers},
  3rd ed. Springer, 2004.
\bibitem{vajda}
  S.~Vajda, \emph{Fibonacci and Lucas Numbers, and the Golden Section}.
  Ellis Horwood, 1989.
\bibitem{koshy}
  T.~Koshy, \emph{Fibonacci and Lucas Numbers with Applications}.
  Wiley, 2001.
\bibitem{bergman}
  G.~Bergman, A number system with an irrational base,
  \emph{Mathematics Magazine} \textbf{31} (1957), 98--110.
\bibitem{coxeter}
  H.~S.~M. Coxeter, \emph{Regular Polytopes}, 3rd ed. Dover, 1973.
\bibitem{grove}
  L.~C. Grove and C.~T. Benson, \emph{Finite Reflection Groups},
  2nd ed. Springer, 1985.
\bibitem{shechtman}
  D.~Shechtman, I.~Blech, D.~Gratias, and J.~W. Cahn,
  Metallic phase with long-range orientational order and no
  translational symmetry,
  \emph{Physical Review Letters} \textbf{53} (1984), 1951--1953.
\bibitem{kramer}
  P.~Kramer and R.~Neri,
  On periodic and non-periodic space fillings of $E^m$ obtained
  by projection,
  \emph{Acta Crystallographica} \textbf{A40} (1984), 580--587.
\bibitem{levine}
  D.~Levine and P.~J. Steinhardt,
  Quasicrystals: a new class of ordered structures,
  \emph{Physical Review Letters} \textbf{53} (1984), 2477--2480.
\bibitem{senechal}
  M.~Senechal, \emph{Quasicrystals and Geometry}.
  Cambridge University Press, 1995.
\bibitem{lewis1975}
  J.~L. Lewis and P.~H. Dworkin,
  Issues in computing with a logarithmic number system,
  \emph{Proc.\ 6th IEEE Symp.\ Computer Arithmetic} (1983), 126--132.
\bibitem{6502}
  R.~Russel, Multiplying with look-up tables on the 6502,
  \emph{Dr.\ Dobb's Journal} (1977).
\bibitem{volder1959}
  J.~E. Volder, The CORDIC trigonometric computing technique,
  \emph{IRE Transactions on Electronic Computers} \textbf{EC-8}
  (1959), 330--334.
\bibitem{metal}
  Apple Inc., \emph{Metal Shading Language Specification}, v3.1.
  Apple Developer Documentation, 2023.
\bibitem{cellular}
  S.~Wolfram, \emph{A New Kind of Science}. Wolfram Media, 2002.
\bibitem{inference_energy}
  International Energy Agency,
  \emph{Electricity 2024 --- Analysis and Forecast to 2026}.
  IEA, Paris, 2024.
  [AI inference accounts for an estimated 80--90\% of lifecycle AI
  energy once models are deployed at scale.]
\end{thebibliography}
\end{document}